\documentclass[10pt,a4paper]{amsart}
\usepackage[margin=19mm]{geometry}
\usepackage{amsmath,amssymb,mathtools}
\usepackage[scr=rsfs,cal=euler]{mathalfa}
\usepackage{xfrac}
\usepackage[unicode,hidelinks,bookmarks=false]{hyperref}
\newcommand{\ie}{i.e.\ }
\newcommand{\cf}{cf.\ }
\newcommand{\resp}{respectively\ }
\newcommand{\Subsection}{\subsection}
\providecommand{\nobreakdash}{\nobreak\hbox{-}}
\setlength{\emergencystretch}{2em}
\multlinegap0pt
\usepackage{bm}
\usepackage{bbm}
\usepackage{dsfont}
\usepackage{xspace}
\usepackage{cancel}
\usepackage{mathtools}
\usepackage{amsthm}
\makeatletter
\@ifundefined{theorem}{\newtheorem{theorem}{Theorem}}{}
\makeatother
\def\e{\mathbf e}
\title[A note on helices in the Euclidean space and in the hyperbol]{A note on helices in the Euclidean space and in the hyperbolic space}
\author{Rafael L\'opez}
\date{Source: arXiv:2507.12861v1 (CC BY 4.0)}
\begin{document}
\maketitle

\noindent\emph{Source and licence.} A note on helices in the Euclidean space and in the hyperbolic space. Version arXiv:2507.12861v1 (CC BY 4.0), distributed under the Creative Commons Attribution 4.0 International licence, as recorded in the arXiv licence metadata for this exact version and shown on the abstract page https://arxiv.org/abs/2507.12861v1. Full rights notice: licenses/arxiv-2507.12861-CC-BY-4.0.txt.\par\smallskip

\noindent\textit{Editorial scope.} Counted unit: the Theorem whose source label is \texttt{t31} in the article (source0.tex lines 221--231, source label \texttt{t31}), delivered together with the article's own complete original proof of it (source0.tex lines 232--269, one \texttt{proof} environment of 38 lines). No mathematics is written, repaired or re-derived: statement and proof are the article's own text line for line. Required context retained uncounted: r (lines 190--192). Every definition, display and label of those lines is retained unchanged. Omitted from this package and not redistributed: 1--189, 193--220, 270--445. Nothing inside a retained excerpt is omitted or altered: every retained body is delivered exactly as the article's lines are written, so no omission receipt of any kind accompanies this package. Citations: the retained excerpts carry no citation command at all, so no bibliography entry of the article is transcribed and no reference is redistributed. Reference adaptation: none was needed and none was made; every reference inside the delivered unit resolves inside it, so no unresolved reference or citation is produced. Printed slips kept literal: no printed word, symbol, hypothesis or number of the counted statement or of its proof was altered. Layout and class: the article's own class is \texttt{amsart} with packages including amssymb, amsthm, graphicx, color, url; the wrapper is the lane's pinned \texttt{amsart} editorial wrapper at 10pt on A4, which restates the theorem-like environments the retained text needs and adds the article's own macro definitions that the retained text needs (\texttt{\textbackslash e}), with extra wrapper packages (bm, bbm, dsfont, xspace, cancel, mathtools). The delivered numbering is the unit's own. Assets: no retained span contains a figure environment, a graphics-inclusion command, a PDF-page inclusion, a table or a TikZ picture; no image file is copied into this package, the package declares no asset dependency and no image file is redistributed with it. Licence: version arXiv:2507.12861v1 of this article is distributed under the Creative Commons Attribution 4.0 International licence, as recorded in the arXiv licence metadata for this exact version and shown on the abstract page https://arxiv.org/abs/2507.12861v1. A full rights notice is delivered at licenses/arxiv-2507.12861-CC-BY-4.0.txt. Characters: every retained excerpt is pure ASCII, so the delivered engine, XeLaTeX with its default Latin Modern text font, reports no missing character.
\begin{equation}\label{r}
\nabla_XY=\nabla^e_XY-\frac{X_3}{z}Y-\frac{Y_3}{z}X+\frac{\langle X,Y\rangle_e}{z}\e_3,
\end{equation}

  \begin{theorem}\label{t31}
   Helices with axis the Killing vector field $V=x\partial_x+y\partial_y+z\partial_z$ and making an angle $\theta$ are the $z$-axis of equation $\{x=0,y=0\}$ or the curves parametrized by 
  \begin{equation}\label{h0}
 \gamma(s)=me^{\frac{\cos\theta}{\sqrt{1+c^2}}s}\left(c \cos\frac{s\sin\theta}{c},c \sin \frac{s\sin\theta}{c}, 1\right).
 \end{equation}
where $m,c>0$. The curvature $\kappa$ and the torsion $\tau$ of $\gamma$   are
\begin{equation}\label{h01}
\kappa=\frac{c^2+\sin^2\theta}{c\sqrt{1+c^2}},\quad \tau=\frac{\sin\theta\cos\theta}{c\sqrt{1+c^2}}.
\end{equation}
Moreover these helices are contained in the Killing cylinder of equation $x^2+y^2=c^2z^2$ and they are geodesics in this cylinder.
  \end{theorem}

  \begin{proof}
    Let $\gamma(s)$, $s\in I$, be a curve such that $V$ is constant along $\gamma$. If $\gamma(s)=(x(s),y(s),z(s))$, then there is a constant $c\geq 0$  such that 
  $$\frac{x^2+y^2+z^2}{z^2}=1+c^2.$$
  If $c=0$, then $\gamma$ is the geodesic given by the vertical line $\{x=0,y=0\}$. This curve makes an angle $\theta=0$ with $V$. Notice that $V(s)={\partial_z}_{|\gamma(s)}$.
  
  From now, we suppose $c>0$. Then  $\gamma$ is contained in the Euclidean  cone $\mathcal{C}_c$ of equation $x^2+y^2=c^2 z^2$. In hyperbolic geometry, the surface $\mathcal{C}_c$ is a Killing cylinder  because it is the locus of points equidistant from the geodesic $\{x=0,y=0\}$.
  
  Since $\gamma$ is contained in $\mathcal{C}_c$,   there exists a smooth function $\varphi=\varphi(s)$ such that 
 \begin{equation}\label{p1}
 \gamma(s)=(cz(s)\cos\varphi(s),cz(s)\sin\varphi(s),z(s)).
 \end{equation}
 Let $\theta$ be the angle between $\gamma$ and $V$. The modulus of $V$ is obtained from 
 $$|V|^2=\langle \gamma,\gamma\rangle=\frac{x^2+y^2+z^2}{z^2}=1+c^2.$$
 Then
 $$\cos\theta=\frac{\langle V,\gamma'\rangle}{|V(\gamma)|}=\frac{xx'+yy'+zz'}{\sqrt{1+c^2}z^2}=\sqrt{1+c^2}\frac{ z'(s)}{z(s)}.$$
 By solving this ODE,  there is $m>0$ such that
 $$z(s)=me^{\frac{\cos\theta}{\sqrt{1+c^2}}s}.$$
 Using that $\gamma$ is parametrized by arc-length, then $x'^2+y'^2+z'^2-z^2=0$. This gives 
 $\varphi'(s)=\frac{\sin\theta}{c}$. Thus, up to a constant on the parameter $s$, we have  $\varphi(s)=\frac{\sin\theta}{c}s$. Definitively, this defines the curve $\gamma$  is given by \eqref{h0}.
 
  We compute the curvature and the torsion of $\gamma$. First, we calculated the normal vector $N$ of $\gamma$. By using \eqref{r}, we have 
\begin{equation}\label{normal}
\begin{split}
\nabla_TT&=\nabla^e_TT-2\frac{T_3}{z}T+\frac{\langle T,T\rangle_e}{z}\e_3\\
&=\frac{m(c^2+\sin^2\theta)}{1+c^2}e^{\frac{\cos\theta}{\sqrt{1+c^2}}s}\left(-\frac{1 }{  c}\cos \left(\frac{s \sin (\theta )}{c}\right),\frac{ - 1}{ c}\sin \left(\frac{s \sin (\theta )}{c}\right) , 1 \right). 
\end{split}
\end{equation}
Computing the modulus, we obtain $\kappa$ given in \eqref{h01}. Hence we have  the unit normal vector by $N=\nabla_TT/k$. Now the computation of the torsion given in \eqref{h01} is obtained by the calculation of  $\nabla_TN$. 

In order to prove that the helices \eqref{h0} are geodesics  in the cylinder $\mathcal{C}_c$ it is enough to check that the normal vector $N$ of $\gamma$ coincides with that unit normal vector $G$ of $\mathcal{C}_c$ along $\gamma$. The cylinder $\mathcal{C}_c$ is given by the implicit equation $x^2+y^2-c^2z^2=0$. Since the hyperbolic metric is conformal to the Euclidean one, the vector $G$ is proportional to the Euclidean gradient $\nabla^e$ of the function $f(x,y,z)= x^2+y^2-c^2z^2$. Then 
\begin{equation*}
\begin{split}
\nabla^ef_{|\gamma}&=2(x,y,-c^2 z)_{|\gamma}\\
&=2me^{\frac{\cos\theta}{\sqrt{1+c^2}}s}\left(c \cos\frac{s\sin\theta}{c},c \sin \frac{s\sin\theta}{c}, -c^2\right)\\
\end{split}
\end{equation*}
This vector is proportional to $\nabla_TT$, \eqref{normal}, and thus it is proportional to $N$.
  \end{proof}

\end{document}
